\scanpage{114}
\mspage{29}
% archive sheet number: 113
4.6. Supposons, \deleted{\unsure{que dans}} dans les conditions de 4.4, que l'on ait pour tout $1\leq i\leq r$

\[
\unsure{D_i|Z_i}=V(a_i),\qquad\text{où }a_i\in\Gamma(\underline{\unsure{O_X|\mathcal O_X}}).
\]

Soit $m$ un entier premier aux car. résiduelles de $X$ \deleted{\unsure{\(\sum D_i\)}} et considérons le schéma

\[
X'=\operatorname{Spec}\unsure{\mathcal O_X|\mathscr O_X}[T_1,\ldots,T_r]/(T_1^m-a_1,\ldots,T_r^m-a_r).
\]

Il est \unsure{alors|donc} bien connu que $X'$ est un schéma régulier, et que \deleted{\unsure{les}} si pour tout $1\leq i\leq r$, \deleted{\unsure{\(a_i\)}} $b_i\in\Gamma(X',\mathcal O_{X'})$ est la section déduite de $T_i$, alors \inserted{les} $\deleted{\unsure{D}}\Delta_i=V(b_i)$ sont des diviseurs réguliers \unsure{de|dans} $X'$ tels que: $\Delta=\sum\Delta_i$ \unsure{soit|est} un diviseur à croisements normaux (cf. SGA 1 $\overline{\underline{\mathrm{XII}}}$ \unsure{\ldots?}) \inserted{Donc on peut appliquer 4.5. à $X'$ et $U'=X'-\operatorname{Supp}\Delta$.} D'ailleurs, il est immédiat que $\operatorname{Supp}\Delta_i$ n'est autre que l'image inverse de $\operatorname{Supp}D_i$, plus précisément, si $p:X'\unsure{{-}|\longrightarrow}X$ est la projection, on a

\[
p^*(D_i)=m\Delta_i.
\]

On en conclut en particulier

\[
U'=p^{-1}(U).
\]

Considérons l'hom. \deleted{\unsure{can.}} de changement de base

\[
p^*(R^*f_*(\Lambda_U))\longrightarrow R^*f'_*(\Lambda_{U'}),
\tag*{(4.\unsure{6|5}.1)}
\]

où $f':U'-X'$ est l'inclusion. \deleted{\unsure{On}} On constate \inserted{\unsure{alors|donc}} aussitôt sur la définition des isomorphismes (4.5.1) que l'on a commutativité dans le diagramme
